\chapter{Background}

\section{Authenticated Encryption with Associated Data (AEAD)}

Cryptographic protocols for embedded communication commonly require three
security objectives simultaneously: confidentiality (to prevent unauthorized
reading of the message), integrity (to detect any modification during
transmission), and authenticity (to verify that the message originates from a
legitimate sender).
Traditionally, these properties were achieved by combining a separate cipher
with a message authentication code (MAC).
However, the secure composition of two independent primitives is non-trivial:
the order of operations matters, and incorrect combinations can introduce
vulnerabilities even when each primitive is individually sound.
Authenticated Encryption with Associated Data (AEAD) addresses this by combining
all three objectives into a single, unified primitive.

An AEAD scheme takes four inputs: a secret key~$K$, a nonce~$N$ (a value that
must be unique per encryption under a given key), a plaintext message~$M$, and a
block of associated data~$AD$.
The associated data is authenticated but not encrypted: it remains readable in
the clear while being bound to the ciphertext via the authentication tag.
The scheme produces a ciphertext~$C$ and an authentication tag~$T$.
Decryption takes the same key and nonce, as well as the ciphertext, tag, and
associated data, and returns the recovered plaintext only if the tag verifies
successfully; otherwise it returns an explicit failure.
This fail-closed behavior is essential in embedded contexts, since a receiver
that processes unauthenticated data before verification completes is vulnerable
to plaintext-dependent side channels and fault injection.

For the constrained microcontroller setting considered in this work, two AEAD
properties are of particular practical relevance.
First, single-pass processing absorbs plaintext and produces the authentication
tag in one sequential pass over the data, minimizing RAM requirements by
avoiding the need to buffer the full message.
Second, a compact state and key schedule reduce code size and initialization
overhead, both of which are essential on devices with limited Flash.
These properties will later serve as evaluation criteria when comparing
AES-128-GCM and Ascon-AEAD128.

\section{AES-128-GCM}

\subsection{The AES Block Cipher}

The Advanced Encryption Standard (AES) is a substitution-permutation block
cipher standardized by NIST in 2001~\cite{FIPS197}.
It operates on a fixed block size of 128 bits and supports key lengths of 128,
192, or 256 bits; this work uses the 128-bit key variant exclusively
(AES-128).
The cipher applies ten rounds of transformation to a $4\times4$ byte state
matrix.
Each round consists of four operations: \texttt{SubBytes} (a non-linear byte
substitution via the AES S-box), \texttt{ShiftRows} (a cyclic shift of state
rows), \texttt{MixColumns} (a linear diffusion step operating on state columns
in $\mathrm{GF}(2^8)$), and \texttt{AddRoundKey} (XOR of the state with a
round-specific subkey).
The 128-bit key is expanded into eleven round keys before encryption begins,
a process known as the key schedule.

The security of AES rests on more than two decades of public cryptanalysis; no
practical attack against the full cipher is known.
With hardware acceleration (a dedicated AES engine), the cipher executes at
throughputs measured in gigabits per second.
On a software-only Cortex-M0 implementation, however, the \texttt{MixColumns}
step and the S-box lookups impose a considerable cycle cost, because the M0
lacks the wider data path and lookup acceleration found on higher-end cores.
This performance gap on minimal instruction-set architectures is the central
motivation for evaluating a lightweight alternative.

\subsection{Galois/Counter Mode (GCM)}

AES alone is a block cipher and not an AEAD scheme: it encrypts exactly one
128-bit block and provides no authentication.
Galois/Counter Mode (GCM) combines AES in counter mode (CTR) for encryption
with a polynomial hash (GHASH) for authentication, yielding an AEAD
construction~\cite{NISTSP80038D}.

In CTR mode, a counter value is encrypted by AES and the resulting keystream
block is XORed with the plaintext.
The counter is incremented for each successive block, which enables
parallelized encryption independent of plaintext content.
GHASH computes a 128-bit authentication tag using polynomial operations over
the ciphertext blocks and associated data in $\mathrm{GF}(2^{128})$, a finite
field defined over 128-bit values.
The hash key~$H$ is derived by encrypting a zero block under the session key
at initialization.

GCM's parallelizability is its primary advantage on general-purpose hardware
with wide data paths.
Its principal weakness on constrained processors is the GHASH step:
multiplication in $\mathrm{GF}(2^{128})$ requires either a hardware carry-less
multiply instruction or a software implementation using multiple 32-bit
operations.
The Cortex-M0 provides only a 32-bit multiply without a carry-less variant, so
a software GHASH implementation on that core requires a series of XOR and shift
operations that contribute noticeably to the total cycle count.
Additionally, GCM requires a unique nonce per encryption; nonce reuse under the
same key is catastrophic, recovering both the hash key and the plaintext with a
small number of observed ciphertexts.

\subsection{Software-Only Scope}

The STM32L476RG does not include a hardware AES peripheral; that block is
present only in the STM32L486 and related parts of the same product family, as
confirmed by the device reference manual.
Both benchmark boards therefore present an identical software-only execution
environment for this comparison.
The absence of a hardware cryptographic accelerator on the selected boards is a
deliberate design choice: restricting the comparison to software implementations
throughout makes the results applicable to the broad class of constrained 32-bit
microcontrollers that ship without vendor-specific AES acceleration, which
represents the majority of cost-sensitive automotive nodes.

\section{The ASCON Family}

\subsection{The ASCON Permutation}

All ASCON variants use the same underlying permutation, written as $p^a$ or
$p^b$ depending on how many rounds are applied.
The permutation operates on a 320-bit state organized as five 64-bit words
$x_0$ through $x_4$.
Each round applies three layers in sequence.

The first layer, $p_C$ (constant addition), XORs a round-specific constant into
word~$x_2$, introducing asymmetry between rounds to prevent slide attacks.
The second layer, $p_S$ (substitution), applies a 5-bit S-box in parallel to
all 64 bit-sliced positions across the five state words.
This S-box provides non-linearity with simple arithmetic and was chosen for the
efficiency of its bitsliced implementation on processors with native 64-bit or
32-bit word operations.
The third layer, $p_L$ (linear diffusion), applies independent linear
transformations to each of the five state words using rotations and XOR,
ensuring that differences propagate rapidly across the full state.

The full ASCON permutation uses twelve rounds ($p^{12}$) for initialization and
finalization.
The internal data-processing step uses eight rounds ($p^8$) for Ascon-AEAD128.
On 32-bit microcontrollers the five 64-bit state words are held in pairs of
32-bit registers, so the rotation-and-XOR structure of the linear layer maps
efficiently onto the available register file without memory traffic.

\subsection{Sponge and Duplex Construction}

ASCON uses a sponge-based construction --- specifically the duplex variant ---
rather than the block-cipher-mode approach used by AES-GCM.
The sponge model divides the state into two parts: the rate~$r$ (the part that
interfaces with input and output) and the capacity~$c$ (the part that remains
hidden from I/O, forming the security margin).
For Ascon-AEAD128, the rate is 128 bits and the capacity is 192 bits, giving a
total state of 320 bits.

In the duplex variant used for AEAD, plaintext is XORed into the rate portion of
the state one rate-sized block at a time, followed by the permutation.
The ciphertext is extracted from the state after the absorption, mixing
encryption and authentication in a single sequential pass.
This single-pass property means that memory buffering of the full message is
never required, a significant advantage over constructions that must process the
ciphertext twice (once for decryption, once for authentication).

The capacity portion of the state is never directly observable by an attacker
and provides the security guarantee.
For Ascon-AEAD128, the 192-bit capacity yields a claimed security level of
128 bits against generic attacks on both confidentiality and integrity.

\subsection{Ascon-AEAD128}

Ascon-AEAD128 is the primary AEAD member of the ASCON family and the algorithm
selected as the NIST Lightweight Cryptography standard~\cite{ASCON2023}.
It accepts a 128-bit key, a 128-bit nonce, associated data of arbitrary length,
and a plaintext of arbitrary length, and produces a ciphertext of the same
length as the plaintext together with a 128-bit authentication tag.

The processing pipeline proceeds in four phases.
In the initialization phase, the 320-bit state is assembled from the key, nonce,
and algorithm-specific initialization vector, then transformed by $p^{12}$.
In the associated data phase, the associated data is absorbed into the rate in
128-bit blocks, each followed by $p^8$; the last block is padded, and a domain
separation constant is XORed into the state at the end of this phase.
In the encryption phase, plaintext blocks are absorbed and ciphertext blocks are
extracted in the same pass, each followed by $p^8$.
In the finalization phase, the key is XORed into the state, $p^{12}$ is applied,
and the authentication tag is extracted from the last 128 bits of the state.

For the benchmark in this work, Ascon-AEAD128 is evaluated using the official
ASCON reference implementation in C~\cite{ASCONSoftware}.

\subsection{Ascon-Hash256}

Ascon-Hash256 is the hash function member of the ASCON family.
It produces a 256-bit digest from an input of variable length using the same
$p^{12}$ permutation as Ascon-AEAD128.
Message absorption proceeds in 64-bit rate blocks, each followed by $p^{12}$,
and the digest is squeezed in two 64-bit output blocks after all input has been
absorbed.


\subsection{Ascon-80pq and the Post-Quantum Outlook}

Ascon-80pq is a variant of Ascon-AEAD128 that extends the key length from
128 to 160 bits while keeping the nonce at 128 bits and the tag at 128 bits.
The longer key is a precautionary measure against Grover's algorithm, which
effectively halves the security level of a symmetric key in bits against an
adversary with a large-scale quantum computer.
A 128-bit key offers approximately 64-bit post-quantum security under Grover's
model, while the 160-bit key of Ascon-80pq provides approximately 80-bit
post-quantum security.

The performance overhead of the 160-bit key variant relative to Ascon-AEAD128
is limited to initialization, where the additional 32 bits are incorporated into
the state setup.
The permutation rounds and data processing phases are identical.
Ascon-80pq is therefore of interest as a forward-compatible migration path for
automotive security architectures that must account for long vehicle service
lifetimes, during which the practical capabilities of quantum computing may
evolve.
Ascon-80pq is part of the original ASCON v1.2 competition submission and is not
included in NIST SP~800-232, which standardizes only Ascon-AEAD128,
Ascon-Hash256, Ascon-XOF128, and Ascon-CXOF128~\cite{ASCONSpec}.

\section{Architectural Comparison: Sponge vs.\ Block Cipher + GCM}

The structural differences between AES-128-GCM and Ascon-AEAD128 have concrete
implications for software performance on constrained microcontrollers, which are
worth making explicit before the benchmark results are presented.

AES-128-GCM is built from two largely independent components: the AES block
cipher and the GHASH authenticator.
Each has its own state, key material, and processing requirements.
The AES key schedule expands 128 bits into 176 bytes of round keys that must
either be precomputed and stored in RAM or recomputed for each execution.
The GHASH state requires an additional 128-bit accumulator and the hash key~$H$.
On a Cortex-M0 without hardware multiply, GHASH is the performance bottleneck
for short messages, since its initialization cost (one AES block encryption) is
spread poorly at small input sizes.

Ascon-AEAD128 holds its entire state in the 320-bit permutation input.
There is no separate key schedule: the key is absorbed directly into the state
during initialization.
The permutation is defined entirely by bitwise rotations and XOR, operations
that are cheap on any 32-bit or 64-bit processor without specialized hardware.
The absence of S-box table lookups also eliminates a class of cache-timing
side channels that must be considered when deploying AES on processors without
constant-time hardware support.

Both algorithms process data in 128-bit units per core operation: AES encrypts
one 128-bit block per cipher call, and Ascon-AEAD128 absorbs one 128-bit rate
block per permutation call.
The throughput difference between them therefore cannot be attributed to a rate
disadvantage on either side and must be found in the per-operation cost of the
underlying primitive.
Whether Ascon-AEAD128 or AES-128-GCM requires fewer cycles on a given platform
cannot be determined theoretically; answering it quantitatively is the central
goal of this benchmark.

\section{Target Platforms}

\subsection{ARM Cortex-M0 --- NUCLEO-F072RB}

The NUCLEO-F072RB board is built around the STM32F072RB microcontroller, which
integrates an ARM Cortex-M0 core clocked at 48\,MHz.
The Cortex-M0 is the smallest and most constrained member of the ARM Cortex-M
family.
It implements the ARMv6-M instruction set architecture, a minimal 16/32-bit
Thumb instruction set without the DSP extensions, floating-point unit, or
hardware divide instruction present in larger M-class cores.
The core has no instruction or data caches.

Relevant characteristics for this benchmark include the absence of a hardware
multiply-accumulate that would accelerate GHASH, a two-cycle latency for the
available $32\times32$-bit multiply (which produces a 32-bit result), and the
lack of a Data Watchpoint and Trace (DWT) unit with a cycle counter.
Because the DWT cycle counter is not present on the M0, cycle-accurate
measurement on this board uses TIM2, configured as a 32-bit free-running counter
driven by the system clock.

The F072RB is representative of the lowest tier of 32-bit microcontrollers
found in automotive body and sensor applications, where unit cost and power
consumption are the dominant design constraints.

\subsection{ARM Cortex-M4F --- NUCLEO-L476RG}

The NUCLEO-L476RG board is built around the STM32L476RG microcontroller, which
integrates an ARM Cortex-M4F core clocked at up to 80\,MHz.
The Cortex-M4F implements the ARMv7E-M instruction set architecture, which
extends ARMv6-M with DSP-oriented instructions including a single-cycle
$32\times32$-bit multiply with 64-bit accumulate (\texttt{MLA},
\texttt{SMLAL}), single-instruction saturating arithmetic, and a
single-precision floating-point unit (FPU).
The core also includes an Adaptive Real-Time (ART) memory accelerator, which
implements an instruction prefetch buffer and a branch cache to reduce Flash
access latency at high operating frequencies.

For timing, the Cortex-M4F provides a DWT unit including the
\texttt{DWT->CYCCNT} cycle counter, a 32-bit register incremented on every
processor clock, readable with a single memory-mapped register access.

The L476RG is representative of a mid-range automotive microcontroller, suitable
for body control modules and sensor fusion nodes with moderate processing
requirements.
Together with the F072RB, the two boards span a performance range relevant to
the lower half of the automotive ECU spectrum.

\section{Related Work}

Performance comparisons between ASCON and AES-GCM on embedded hardware have
appeared in several contexts since ASCON's selection as the NIST LWC standard.
Benchmarks conducted on AVR 8-bit and MSP430 16-bit platforms during the NIST
LWC competition evaluation process demonstrated ASCON's advantage on
ultra-constrained targets, but these results are not directly applicable to
32-bit ARM Cortex-M deployments that represent the current generation of
automotive microcontrollers.

Work targeting ARM Cortex-M platforms specifically has been published both by
the ASCON team and by independent researchers.
The ASCON team provides cycle counts for their reference and optimized
implementations on several Cortex-M targets~\cite{ASCONSoftware}.
However, published figures often report single-point measurements at a fixed
message length rather than throughput curves across a logarithmic range of input
sizes, which are needed to characterize behavior for the short messages typical
in automotive CAN frames and sensor update packets.

To the authors' knowledge, no published work directly compares AES-128-GCM and
Ascon-AEAD128 across both a Cortex-M0 and a Cortex-M4F target with a
controlled, software-only methodology that includes compiler optimization level
as an explicit variable.
This work aims to fill that gap by providing reproducible benchmark data for
both platforms under a unified measurement framework.
